\begin{afterword}
\summaryhead

The author once thought perseverance meant working 14-hour days, which the author cannot do. Looking back on years of work in AI, the author found that almost every important problem had been overestimated. The author first pictured AI as needing forty years and a Manhattan Project, but stayed, and the path came to look shorter. Friendly AI and a precise understanding of intelligence, once written off as impossible, were later taken on. Single problems came to seem smaller while the whole task grew.

``Impossible,'' the post argues, usually means ``I can't see any way to do that,'' a sign of confusion, and confusion is in the map. A problem that seems impossible is not tried, so the confusion does not clear. Most AGI researchers, the post says, expect an easy solution in five years instead. Citing Richard Hamming, it says the most important problems often look impossible, and warns that attempting them is for very special occasions: talent is only the ante, the stake is years of your life, and you can lose. Perseverance works at three timescales: not running away (seconds), working (hours), and sticking at it (years).

\responsehead

The post is memoir and advice, and both mostly hold up. Where the memoir can be checked, it matches the record. The 1996 essay ``Staring into the Singularity'' gives 2035, ``less than forty years!'' ``Coding a Transhuman AI'' was announced in September 1998, a month after Hypermart's sale was reported. ``The Plan to Singularity'' of 1999 set out the open-source plan and aimed at a seed AI by 2010, which is the shortened path the post describes. The advice states its own limits. It tells readers that ``you can lose,'' that most people should ``stick to the possible,'' and that knowing when to lose hope is a skill. It names the ways the attempt fails: blind alleys, and too little ability.

The opening claim is stronger than its evidence. The author ``had overestimated the difficulty of almost every single important problem.'' What the post shows is that the problems came to seem less intimidating. It reports none solved, and it grants that ``you don't know how much work will be left when the confusion clears.'' The post's later sentences speak of ``apparent'' difficulty, and they are the accurate version.

The claims about other researchers rest on the author's impressions. ``By and large,'' most AGI researchers are said to be in the field because they think they have found ``the Key to AI'' that will work ``in just five years.'' Others are said to call solvable-looking problems ``Impossible!'' No case is named. The links go to a Star Wars parody and to a post, not in the book, built on impressions from conferences.

The post cites Richard Hamming's questions about important problems and then says that the most important problems are often impossible. Hamming's talk draws the line differently: what makes a problem important ``is that you have a reasonable attack.'' A problem where one can see no way in is, on his account, not yet an important one. The two agree that drive matters; Hamming says ``The steady application of effort with a little bit more work, intelligently applied is what does it.'' They differ on which problems deserve it.

In short: a memoir that the record bears out and advice that states its own limits, whose opening claim of overestimated difficulty rests on how the problems came to seem, and which pairs Hamming's questions with a view of important problems that his test of ``a reasonable attack'' does not share.
\end{afterword}
